\section{Balanced rounding and realization in the metric}\label{sec:rounding}

Balanced subtree rounding with constant movement loss on separated trees
was developed in \cite[Section~5.2]{BBMN15}. The formulation here also
allows internal parks and a prescribed parent total correlated with the
old child counts.

The allocation theorem allows mass to stop at internal vertices.  We round
both subtree amounts and parks, working from the root toward the leaves.
At each vertex the already rounded parent fixes an integer total for its
children and its park.  That total may be correlated with their old counts.
Our local kernel first meets the total in every realized state, then repairs
the coordinate means by exchanges that preserve each state's total.

The resulting integer parks can be matched along the tree.  Parks matched
at different vertices are paid by tree movement; parks matched at the same
vertex require the separate anchor-travel estimate when their representative
has moved.

Throughout this section, the filtration $\mathcal F_t$ contains the request
prefix and the independent partition tape, as in Section~\ref{sec:allocation}.
It excludes a further independent tape used only for rounding.  All fractional
states and anchors at time $t$ are $\mathcal F_t$-measurable.  Write
\[
 p_v^\circ(t)=q_v(t)-\sum_{i\text{ child of }v}q_i(t),
\]
where the sum is empty at a leaf.  We seek integer subtree counts $n_v(t)$ and
integer parks $c_v(t)$ satisfying
\begin{equation}\label{eq:rounding-invariants}
 \begin{aligned}
 n_{\rm root}(t)&=k,&
 n_v(t)&=c_v(t)+\sum_{i\text{ child of }v}n_i(t),\\
 n_v(t)&\in\{\lfloor q_v(t)\rfloor,\lceil q_v(t)\rceil\},&
 \E[n_v(t)\mid\mathcal F_t]&=q_v(t),\\
 c_v(t)&\in\{\lfloor p_v^\circ(t)\rfloor,
                    \lceil p_v^\circ(t)\rceil\},&
 \E[c_v(t)\mid\mathcal F_t]&=p_v^\circ(t).
 \end{aligned}
\end{equation}
These requirements concern conditional means given $\mathcal F_t$ alone.
We do not prescribe the same means after additionally conditioning on the
realized rounded ancestors; such requirements would generally be inconsistent
with their prescribed integer totals.

\subsection{A finite rounding kernel for one star}

\begin{lemma}[Balanced rounding with a prescribed random total]
\label{lem:rounding-star}
Let $a,b\in[0,\infty)^d$, and let $X\in\Z_{\ge0}^d$ be a random vector with
\[
 X_i\in\{\lfloor a_i\rfloor,\lceil a_i\rceil\},
 \qquad \E X_i=a_i.
\]
Put $M=\sum_iX_i$.  Let $N$ be any random integer, arbitrarily correlated with
$X$, such that
\[
 N\in\left\{\left\lfloor\sum_i b_i\right\rfloor,
                 \left\lceil\sum_i b_i\right\rceil\right\},
 \qquad \E N=\sum_i b_i.
\]
There is a randomized kernel from the finite set of states $(X,N)$ to vectors
$X'$ such that, in every state and every outcome of the kernel,
\[
 X'_i\in\{\lfloor b_i\rfloor,\lceil b_i\rceil\},
 \qquad \sum_iX'_i=N.
\]
Its coordinate means, averaged over both the original state and the kernel,
satisfy
\[
                         \E X'_i=b_i.
\]
With
\[
 T=\sum_i|a_i-b_i|,
 \qquad I=\E|N-M|,
\]
the kernel can be chosen to satisfy
\begin{equation}\label{eq:rounding-star-bound}
             \E\|X'-X\|_1\le32T+11I.
\end{equation}
Moreover, for any balanced random total $N$ with the displayed mean, a balanced
vector with total $N$ and coordinate means $b_i$ exists without specifying an
old vector $X$.
\end{lemma}

\begin{proof}
All expectations in this proof are taken under the given joint distribution;
no independence assumption is made.  Set $l_i=\lfloor b_i\rfloor$ and
$u_i=\lceil b_i\rceil$.  Clip $X_i$ into $[l_i,u_i]$, obtaining $Y_i$.
For each coordinate,
\begin{equation}\label{eq:rounding-clip}
                       \E|Y_i-X_i|\le |a_i-b_i|.
\end{equation}
Indeed, if $a_i\le b_i$, upper clipping is impossible.  Lower clipping either
does nothing, or $l_i>\lfloor a_i\rfloor$, in which case every possible $X_i$
is at most $l_i$ and its expected cost is $l_i-a_i\le b_i-a_i$.
If $a_i\ge b_i$, the symmetric argument applies to upper clipping.
This includes integral targets and equality of endpoints.  Summing
\eqref{eq:rounding-clip} gives $\E\|Y-X\|_1\le T$.

For each state, adjust $Y$ within the integer intervals $[l_i,u_i]$ until its
sum equals the actual value of $N$, obtaining $Z$.  The adjustment is always
possible: the sums of allowed vectors comprise every integer between
$\sum_i l_i$ and $\sum_i u_i$, and both endpoint integers of $\sum_i b_i$
lie in that interval.  Increasing coordinates only, or decreasing coordinates
only, as needed, costs exactly $|N-\sum_iY_i|$.  A fixed coordinate order
makes this choice deterministic.  Thus, writing $L=\E\|Z-X\|_1$,
\begin{equation}\label{eq:rounding-preliminary}
 \begin{aligned}
 L&\le \E\|Y-X\|_1+\E\left|N-\sum_iY_i\right|\\
  &\le 2T+I,\\
 D:=\sum_i|\E Z_i-b_i|
  &\le \sum_i|\E Z_i-a_i|+T
   \le L+T\le3T+I.
 \end{aligned}
\end{equation}
Only the means of $Z$ may now be wrong.  Requiring
$\E[X'_i\mid N]=b_i$ would force $N=\sum_i b_i$ on every
positive-probability event.  We instead repair the means under the given
joint law while preserving each state's actual total.

To construct this repair by finite optimization, let
\[
 H=\{i:b_i\notin\Z\},\qquad f_i=b_i-l_i\quad(i\in H),
 \qquad s=\left\lfloor\sum_{i\in H}f_i\right\rfloor.
\]
An allowed vector is specified by a subset $A\subseteq H$ of coordinates
rounded upward.  Its cardinality is $N-\sum_i l_i$, and hence is either $s$
or $s+1$.  If $\sum_{i\in H}f_i$ is integral, only $s$ occurs.  The preliminary
vector $Z$ specifies such a subset $A_0$ in each original state.

For every original state $\omega=(X,N)$ of positive probability, consider all
probability distributions on subsets of the required cardinality for that
state.  Their product is a nonempty compact finite-dimensional set.  A choice
in this product is precisely a randomized kernel that keeps the observed
total $N$ fixed.  Minimize on it
\begin{equation}\label{eq:rounding-repair-objective}
 \E|A\mathbin{\triangle}A_0|
       +10\sum_{i\in H}|\Pr(i\in A)-f_i|.
\end{equation}
Here $|A\mathbin{\triangle}A_0|$ is the number of coordinates changed; it
equals the $\ell_1$ change of the associated integer vectors.  A minimum
exists by continuity.  Let $p_i=\Pr(i\in A)$ at a minimizer.  Since every
outcome has the required total,
\[
                 \sum_{i\in H}p_i=\E N-\sum_i l_i
                                      =\sum_{i\in H}f_i.
\]
Consequently, if a discrepancy remains, there are distinct indices $i,j$
with $p_i>f_i$ and $p_j<f_j$.

First suppose that a supported state-and-subset pair has $i\in A$, $j\notin A$.
On a sufficiently small amount $\varepsilon>0$ of its probability mass replace
$A$ by $A\setminus\{i\}\cup\{j\}$.  This preserves that state's total.  Choose
$\varepsilon$ no larger than the two discrepancies and the mass of this
state-and-subset pair.  The sum of marginal discrepancies falls by
$2\varepsilon$, whereas the first term in
\eqref{eq:rounding-repair-objective} rises by at most $2\varepsilon$.
This contradicts minimality.

In the remaining case, no supported subset contains $i$ without $j$.
Because $p_i>f_i>0$, some supported subset $A$ contains both.  Because
$p_j<f_j<1$, some supported subset $B$ contains neither.  Their cardinalities
differ by at most one.  There must therefore be an index $h\in B\setminus A$:
otherwise $B\subseteq A\setminus\{i,j\}$ would give
$|B|\le|A|-2$.  On equal sufficiently small probability masses
$\varepsilon$ of these two state-and-subset pairs, make the changes
\[
 A\longmapsto A\setminus\{i\}\cup\{h\},
 \qquad
 B\longmapsto B\setminus\{h\}\cup\{j\}.
\]
Each change occurs within its own original state and preserves that state's
total.  The marginal of $h$ does not change, the marginal of $i$ decreases by
$\varepsilon$, and that of $j$ increases by $\varepsilon$.  The Hamming cost
in \eqref{eq:rounding-repair-objective} rises by at most $4\varepsilon$, by
the triangle inequality for symmetric difference.  The discrepancy penalty
falls by $20\varepsilon$.  Again this is a strict improvement.  This argument
allows the two subsets to belong to different original states and different
values of the rounded parent; it does not alter their probabilities or totals.

Thus every minimizer has $p_i=f_i$ for all $i\in H$.  The other coordinates
are fixed integers, so the corresponding vector $X'$ has all required means.
Comparison with the kernel that always keeps $A_0$ gives
\[
                       \E\|X'-Z\|_1\le10D.
\]
Combining with \eqref{eq:rounding-preliminary} proves
\[
 \E\|X'-X\|_1\le L+10D
       \le(2T+I)+10(3T+I)=32T+11I.
\]

For initialization, choose in each state $N$ any allowed subset of the
required size, and minimize the same objective relative to that subset.
The exchange argument still proves that the minimum has the desired means;
it used only the correct mean and the two possible values of $N$.
If $H$ is empty, $N=\sum_i b_i$ and the sole allowed vector is $b$ itself,
so no repair is needed.  This proves the initialization assertion as well.
\end{proof}

\subsection{The joint rounding process}

We apply Lemma~\ref{lem:rounding-star} in a fixed order from the root toward
the leaves.  At an internal vertex $v$, the coordinates of its star are its
ordinary children together with one extra coordinate for its park.  The
target vector is $(q_i)_{i\text{ child of }v}$ together with $p_v^\circ$,
whose sum is $q_v$.  The prescribed integer total is the already generated
$n_v$.  At a leaf simply set $c_v=n_v$.

At time zero, put $n_{\rm root}=k$ and use the initialization assertion
successively at every internal vertex.  The resulting conditional means are
those in \eqref{eq:rounding-invariants}, by induction over the fixed order.
Consistency holds in every outcome, so summing all parks gives exactly $k$.
No movement is charged for this choice of the initial simulated configuration;
its mismatch from the actual starting tuple is charged once below.

Consider an update from $t-1$ to $t$.  The old rounding variables are functions
of $\mathcal F_{t-1}$ and the independent rounding tape.  The fractional
algorithm and the requests do not depend on that tape.  Therefore revealing
$\mathcal F_t$ preserves their conditional distribution from
$\mathcal F_{t-1}$; in particular,
\begin{equation}\label{eq:old-rounding-means}
 \E[n_v(t-1)\mid\mathcal F_t]=q_v(t-1),
 \qquad
 \E[c_v(t-1)\mid\mathcal F_t]=p_v^\circ(t-1).
\end{equation}
Formally, if $U$ denotes the independent rounding tape and
$Z=g(\mathcal F_{t-1},U)$ is one of these variables, its conditional law given
$\mathcal F_t$ is obtained by integrating $g$ over the unchanged law of $U$.
Its dependence on the observed data is already fixed by
$\mathcal F_{t-1}\subseteq\mathcal F_t$.  This also shows why the additional
rounding tape must not affect any of the fractional observations.

Condition now on a fractional history through time $t$.  At a star $v$, use
the joint law of its old child and parking counts and its newly generated
parent count $n_v(t)$.  This is a finite law determined by the known past
kernels and the kernels already used above $v$ at this step.  The parent
count has conditional mean $q_v(t)$ by the top-down induction, while
\eqref{eq:old-rounding-means} supplies the old coordinate means.  The lemma
therefore applies even though this newly generated parent can be correlated
with all the old local counts.  Select a minimizing finite kernel and sample
it at the actual state encountered.  Subsequent stars do not change counts
already assigned, so the induction gives all of
\eqref{eq:rounding-invariants} jointly.

This description defines causal kernels, rather than merely unrelated
couplings at separate times.  Each kernel uses the joint law conditional on
the fractional history, then acts on the realized local state and parent
total.  It need not impose its marginal targets separately on the events
specifying individual ancestor roundings.  At a fixed finite horizon there
are only finitely many relevant request and partition outcomes: on $Y$, a
radius is needed only through the finite ball it defines, and only finitely
many insertion and retirement draws are used.  We can therefore choose one
minimizing kernel at each relevant history without a measurable-selection
issue.  Choices at histories of probability zero can be arbitrary balanced
continuations.  The distributional policy is allowed to depend on the given
law and horizon at this stage.

\begin{proposition}[Rounding and metric realization]\label{prop:rounding}
Assume the fractional allocations and anchors constructed in the preceding
sections, with $\tau\ge22$.  They admit a causal rounding satisfying
\eqref{eq:rounding-invariants} for which
\begin{equation}\label{eq:rounding-global}
 N_{\rm rnd}:=\E\sum_{t=1}^H\sum_{v\ne{\rm root}}
                 r_{|v|-1}|n_v(t)-n_v(t-1)|\le132V.
\end{equation}
Placing the integer parks at their current anchors gives simulated
configurations $S_t$ of exactly $k$ members containing each current request.
There is a matching of consecutive configurations with expected total cost
at most
\begin{equation}\label{eq:rounding-realization}
                         120N_{\rm rnd}+40T_{\rm anchor}.
\end{equation}
Consequently their expected movement is at most
$C(V+\ell P+\ell Q)+D$ for an absolute $C$ and a finite additive term $D$
depending only on the fixed finite structure.
\end{proposition}

\begin{proof}
We first prove the movement estimate in the tree.  For an internal vertex
$v$ and time $t$, define
\[
 T_{v,t}=\sum_{i\text{ child of }v}|\Delta q_i(t)|
                                      +|\Delta p_v^\circ(t)|,
 \qquad
 I_{v,t}=\E[|\Delta n_v(t)|\mid\mathcal F_t].
\]
The local lemma, conditional on $\mathcal F_t$, gives
\[
 \E\left[\sum_{i\text{ child of }v}|\Delta n_i(t)|
                          +|\Delta c_v(t)|\,\middle|\,\mathcal F_t\right]
                                  \le32T_{v,t}+11I_{v,t}.
\]
Since $p_v^\circ=q_v-\sum_iq_i$,
\[
 T_{v,t}\le2\sum_{i\text{ child of }v}|\Delta q_i(t)|
                                          +|\Delta q_v(t)|.
\]
Multiply by $r_{|v|}$, sum over internal vertices and times, and take
expectations.  Every nonroot vertex occurs once as an ordinary child, and
$\Delta q_{\rm root}=0$.  Thus
\begin{equation}\label{eq:rounding-target-sum}
 \E\sum_{t,v\text{ internal}}r_{|v|}T_{v,t}
                                            \le(2+\tau^{-1})V.
\end{equation}
The same weighting of the parent terms gives
\begin{equation}\label{eq:rounding-parent-sum}
 \E\sum_{t,v\text{ internal}}r_{|v|}I_{v,t}
                                       \le\tau^{-1}N_{\rm rnd},
\end{equation}
because the root count is fixed and the internal nonroot vertices are a
subset of all nonroot vertices.  Dropping the nonnegative parking movements
from the local left-hand sides now yields the exact recurrence
\begin{equation}\label{eq:rounding-absorption}
 N_{\rm rnd}\le(64+32/\tau)V+(11/\tau)N_{\rm rnd}.
\end{equation}
All sums are finite.  For $\tau\ge22$ this gives
$N_{\rm rnd}\le(128+64/\tau)V\le132V$, proving
\eqref{eq:rounding-global}.  In particular, initialization of the rounding
does not introduce a new additive variation term.  The same absolute $\tau$
can satisfy this threshold and the allocation theorem's absorption threshold.

We next describe the matching in the tree.  Give the edge entering a nonroot
vertex $v$ length $r_{|v|-1}$.  The consistency identities in
\eqref{eq:rounding-invariants} imply that $n_v(t)$ is the total number of
integer parks in the subtree at $v$.  First match
$\min\{c_v(t-1),c_v(t)\}$ old and new copies at every identical vertex $v$.
The unmatched copies define integer supplies and demands with equal total.
Starting at the leaves, cancel supplies against demands whenever they meet,
and carry the uncancelled copies to the parent.  At the edge entering $v$,
the number of copies carried is exactly $|\Delta n_v(t)|$: the signed
imbalance in that subtree is $\Delta n_v(t)$, and after cancellation all
remaining copies have the same sign.  The root has zero final imbalance.
Following each matched pair's path therefore gives total tree length exactly
\begin{equation}\label{eq:rounding-transport-cost}
                   \sum_{v\ne{\rm root}}r_{|v|-1}|\Delta n_v(t)|.
\end{equation}
This also proves the lower bound for any matching, since each subtree's net
imbalance must cross its entering edge.  Fixing orders for cancellations
makes the chosen matching a function of the current and previous rounded
states.

Every occupied vertex has a current validity witness.  Indeed $c_v(t)>0$
implies $p_v^\circ(t)>0$, since the only balanced rounding of zero is zero.
Hence $q_v(t)>0$, and the allocation domination $q_v(t)\le C_0m_v(t)$ implies
$m_v(t)>0$.  Some point of $Y$ must then have the prefix defining $v$.
This is also the witness assertion of Lemma~\ref{lem:park-variation}.
At a nonroot vertex of depth $j$, its witness is within $20r_j$ of the
vertex's current anchor, by Lemma~\ref{lem:fixed-alphabet}.

Consider a matched old vertex $u$ and new vertex $v$ with $u\ne v$.
Let their last common prefix be $w$ of depth $j$.  If $j\ge1$, choose an
old witness $y_-$ for $u$ and a new witness $y_+$ for $v$.  Write $a_-(z)$
and $a_+(z)$ for an anchor before and after the step, respectively.
Both witnesses use the same depth-$j$ key at $w$ in the two consecutive
maps.  By the fixed-alphabet construction, its two anchors are either equal
or related by an explicit heavy-ball reuse, and in either case
\[
 d(a_-(w),a_+(w))\le40r_j.
\]
The witnesses are within $20r_j$ of the corresponding anchors at $w$.
They are also within $20r_{|u|}$ and $20r_{|v|}$ of the anchors of their own
endpoints.  Since both endpoint radii are at most $r_j$, the triangle
inequality gives
\begin{equation}\label{eq:rounding-distinct-anchors}
 \begin{aligned}
 d(a_-(u),a_+(v))
 &\le20r_{|u|}+20r_j+40r_j+20r_j+20r_{|v|}\\
 &\le120r_j.
 \end{aligned}
\end{equation}
The tree path between distinct endpoints passes through an edge of length
$r_j$: this remains true when $w=u$ or $w=v$, since then the path starts
with the edge to the next depth.  Thus \eqref{eq:rounding-distinct-anchors}
is at most $120$ times this path's tree length.  If $j=0$, both anchors
belong to $Y$, so their distance is at most $\diam(Y)=r_0$.  Their path
crosses an edge of length $r_0$, also when one endpoint is the root.  The
same bound follows.  Only occupied vertices were used in this reasoning;
no assertion about distances associated with nonexistent prefixes is needed.

It remains to pay for the copies matched at one and the same vertex.
The root anchor is fixed, so it contributes zero.  For a nonroot vertex
$v$ of depth $j$, put
$h_v=\min\{c_v(t-1),c_v(t)\}$.  If its anchor moves while both copies
are occupied, its last label is retained across the step.  Only heavy
labels can move in this way, and the displacement is at most $40r_j$.
Write $E_{t,j,\lambda}$ for the event that label $\lambda$ is retained and
moves at this step.  This event and both anchors are
$\mathcal F_t$-measurable.  Since $h_v\le c_v(t)$, regardless of any
correlations created by the matching,
\[
 \E[h_v\one_{E_{t,j,\lambda}}\mid\mathcal F_t]
       \le p_v^\circ(t)\one_{E_{t,j,\lambda}}.
\]
Summing first over all depth-$j$ vertices with last label $\lambda$ turns
the right-hand side into
$b_{t,j,\lambda}\one_{E_{t,j,\lambda}}$, in the notation of
Section~\ref{sec:travel}.  Summing over labels, levels and times pays the
expected cost of all these coincident-vertex matches by $40T_{\rm anchor}$.
Together with \eqref{eq:rounding-transport-cost} and the bound for distinct
vertices, this proves \eqref{eq:rounding-realization}.

Finally let $x$ be the current request.  Its assigned leaf has $q_v\ge1$
by Theorem~\ref{thm:allocation}, so balancedness gives $c_v=n_v\ge1$.
At level $J$, a point assigned to a ball key is within $20r_J$ of that
key's center.  Since $40r_J$ is smaller than the least positive distance
in $Y$, this point equals the center.  A fallback key is anchored at its
point as well.  The anchor of the request leaf is therefore exactly $x$,
and the simulated configuration serves it.  The claimed final estimate now
follows from \eqref{eq:rounding-global}, \eqref{eq:rounding-realization},
and Proposition~\ref{prop:anchor-travel}.
\end{proof}

\subsection{Uniform finite-law synthesis}

We can now complete the distributional statement in the setup.  We give the
additive accounting explicitly because its independence of the horizon is
essential for the transfer in the next section.

\begin{proof}[Proof of Proposition~\ref{prop:finite-law}]
Fix the finite metric $Y$, the server number $k$ and the initial tuple $s$.
Choose the absolute separation factor once, sufficiently large for
Theorem~\ref{thm:allocation} and with $\tau\ge22$.  Fix the scales and the
finite alphabets from Lemma~\ref{lem:fixed-alphabet}.  These choices depend
only on $Y,k$ and the absolute construction parameters.  For a given
finite-support input law at a finite horizon, choose its hidden optimal
trajectory and carry out the partition, allocation, anchor and rounding
constructions.  Their interfaces give absolute constants $C_P,C_A,C_T$ and
finite remainders $D_P,D_A,D_T$ such that
\begin{equation}\label{eq:rounding-budget-summary}
 \begin{aligned}
 P&\le C_P\ell Q+D_P,
       &&\text{by Proposition~\ref{prop:embedded-cost}},\\
 V&\le C_A\ell P+D_A,
       &&\text{by Theorem~\ref{thm:allocation}},\\
 T_{\rm anchor}&\le C_T(V+\ell P+\ell Q)+D_T,
       &&\text{by Proposition~\ref{prop:anchor-travel}}.
 \end{aligned}
\end{equation}
In particular the conditional laws used by the allocation are those of
the hidden members under the final online map.  Filtering keeps the old
map with the old hidden variable, and drift uses the coupled old and new
maps, exactly as required by the first two estimates.  Rounding neither
changes those observations nor supplies information to the partition rules.

Here is a census of the fixed objects and their possible endpoint charges
in \eqref{eq:rounding-budget-summary}.  Let $\mathcal T$ be the full finite
prefix tree, let $d_v$ be the number of children of an internal vertex,
and let $L_j$ be the number of label slots at level $j$.  These numbers can
be large.  They are fixed before the law and horizon are chosen.

\begin{enumerate}
\item At each fixed vertex and rank, the elementary tracking potential is
$p(1-p)+(p-p^\circ)^2\in[0,1]$.  There are finitely many such rank,
flag and endpoint-function trackers, indexed by the $k$ ranks of the
fixed vertices, with the fixed finite number of versions used in the
allocation construction.  Their tolerances, including the tolerances of
order $1/\ell$, have already been fixed.  The resulting initialization and
endpoint charges are finite functions of $k$ and the vertex index, rather
than of a trajectory's number of resets.

At an internal vertex, the sizes, core sums and cap scales are bounded by
fixed multiples of $k$.  The proportion variables have the fixed simplex
or bounded-$\ell_1$ domains used in their definitions.  The primitive bounds
in Section~\ref{sec:proportions} therefore bound the absolute values of
both proportion potentials by a finite quantity, for example
$C\ell k(d_v+1)$ after increasing an absolute constant.  The switching
potential is bounded by $Ck$.  Initial size and epoch charges likewise
involve only these bounded sizes.  Multiplying these finitely many bounds
by their fixed star weights $r_{|v|}$ and by the fixed coefficients in the
local estimates gives a finite common choice of $D_A$.

\item At scale $r_j$, each pilot potential lies in $[0,2r_jk]$.
There is one heavy template and a finite number of tier templates, the
latter with their fixed weights $e^{-h}$.  The heavy distance-ramp and the
tier noncoverage potentials lie in $[0,r_jk]$ for each template.  Their
initial and terminal contributions are bounded by the corresponding
finite sums of $2r_jk$ and $r_jk$ over levels and tiers, with the fixed
coefficients from Section~\ref{sec:partitions}.  Hence $D_P$ can be
chosen uniformly.

\item Each held key size in Section~\ref{sec:travel} is in $[1,k+1]$.
The affine tangent potential to $-s\log s$, evaluated at an argument and
reference in that interval, has absolute value at most
$C(k+1)\ell$.  The normalized travel potential for a level-$j$ label is
nonnegative and at most $Cr_jk$, because its parked-mass coefficient
$b/\widehat m$ is bounded by an absolute constant and its counting
integral is at most $k$.  The sums
\[
 \sum_{j=1}^J L_j r_j(k+1)\ell
       \quad\hbox{and}\quad
 \sum_{j=1}^J L_j r_jk
\]
therefore bound all relevant endpoint terms up to fixed coefficients.
They give a finite common choice of $D_T$, including its size-tracker
initialization charges.

\item The rounding process has no additional variation remainder, by
\eqref{eq:rounding-global}.  Its initial configuration $S_0$ is a multiset
of $k$ points in $Y$.  Consequently the initial minimum matching distance
from $s$ to $S_0$ is at most $k\diam(Y)$, for every rounding outcome.
\end{enumerate}

This list uses persistent state at every fixed slot.  A rank or key tracker
continues through absence and later reappearance of its label; no new
initialization allowance is charged when a slot is reused.  Epoch changes
and normalization refreshes are paid transitions in the preceding
estimates.  The alphabets have bounded size because each tier and scale
has at most $|Y|K^2$ relevant roster entries, with fixed additional slots
for the label exclusion rule.  Fresh random draws at successive insertions
do not create fresh tree vertices or fresh potentials.  Thus the finite
endpoint bounds above hold uniformly over all possible histories, and
are summed over slots and templates once.  The number of histories, the
number of label reuses and the horizon do not enter $D_P,D_A,D_T$.

To make the final combination explicit, put
\[
 C_1=120\cdot132+40C_T,
 \qquad C_2=40C_T,
 \qquad C_*=C_1C_A+C_2.
\]
Proposition~\ref{prop:rounding} and \eqref{eq:rounding-budget-summary} give
\begin{align*}
 \E[\text{simulated movement}]
 &\le C_1V+C_2\ell P+C_2\ell Q+40D_T\\
 &\le C_*\ell P+C_2\ell Q+C_1D_A+40D_T\\
 &\le (C_*C_P+C_2)\ell^2Q
                +C_*\ell D_P+C_1D_A+40D_T,
\end{align*}
where $\ell\ge1$ was used in the last line.  Apply
Lemma~\ref{lem:lazy-simulation} from the actual labeled tuple $s$.
It yields a causal strongly lazy policy serving the same requests, with
expected cost at most this simulated movement plus the initial matching
mismatch.  Thus Proposition~\ref{prop:finite-law} holds with the absolute
constant
\[
                  a=C_*C_P+C_2
\]
and, for example, the finite nonnegative constant
\begin{equation}\label{eq:rounding-finite-additive}
 D_0=C_*\ell D_P+C_1D_A+40D_T+k\diam(Y).
\end{equation}
All three remainders can be enlarged to nonnegative uniform bounds from the
census above.  Formula~\eqref{eq:rounding-finite-additive} consequently
depends only on the fixed metric, $k$, the start and the fixed construction;
it is independent of the input law and horizon.  Repeated starting
positions cause no change, since both simulated configurations and the
matching simulation treat repeated points as separate copies.
\end{proof}
